\documentclass[11pt,a4paper]{article}
\usepackage[margin=1in]{geometry}
\usepackage{amsmath,amssymb,amsthm}
\usepackage{booktabs}
\usepackage{hyperref}
\usepackage[T1]{fontenc}
\usepackage{lmodern}
\usepackage[htt]{hyphenat}
\usepackage{microtype}
\setlength{\emergencystretch}{3em}
\newtheorem{theorem}{Theorem}
\newtheorem{proposition}[theorem]{Proposition}
\theoremstyle{definition}
\newtheorem{definition}[theorem]{Definition}
\hypersetup{colorlinks=true,linkcolor=blue,urlcolor=blue,citecolor=blue}

\title{\textbf{Uniruled Fourfolds in Lean 4: A Conditional Cycle-Class Reduction under Four Named Hypotheses (L11), (BS), (CL), (U)}}
\author{Benjamin Stanley Frohman\\
\small ORCID: \href{https://orcid.org/0009-0006-7068-3718}{0009-0006-7068-3718} \quad|\quad
GitHub: \href{https://github.com/BenFrohman/uniruled-fourfolds}{\texttt{BenFrohman/uniruled-fourfolds}}\\
\small Zenodo concept DOI (all versions): \href{https://doi.org/10.5281/zenodo.23652973}{\texttt{10.5281/zenodo.23652973}}}
\date{Earliest Commit Provenance: 2026-10-01 04:51:29 PM CDT (UTC-5)\\
Siloed Archival Edition: October 25, 2026 (America/Chicago CDT)\\
Revised: November 23, 2026 (America/Chicago CST), conditional restatement}

\begin{document}
\maketitle

\begin{abstract}
A Lean~4 development (\texttt{UniruledFourfolds/\allowbreak Blueprint.lean}, \texttt{UniruledFourfolds/\allowbreak NonVacuity.lean}, \texttt{UniruledFourfolds/\allowbreak Semantics.lean}, \texttt{UniruledFourfolds/\allowbreak MukaiLattice.lean}) of a conditional cycle-class reduction for smooth complex projective fourfolds. The development contains no \texttt{sorry} and no \texttt{axiom}. It proves: if a correspondence from a threefold satisfies a Bloch--Srinivas type decomposition on Hodge classes and is compatible with the cycle class map, and the Lefschetz theorem on $(1,1)$-classes holds on the threefold, then the codimension-2 cycle class map is surjective onto rational Hodge classes. Each of these inputs is a named hypothesis, not a theorem of the development. Finite models show that the hypotheses are jointly consistent and that none of them can be dropped.
\end{abstract}

\section{Scope, Scientific Honesty, and Non-Claim Boundary}
Every result in this paper is conditional on the named hypotheses (L11), (BS), (CL) and (U) of Section~\ref{sec:setting}. None of these hypotheses is proved in this repository. The groups, cycle class maps and the uniruledness predicate are uninterpreted data in Lean: the development does not construct cohomology, Chow groups or cycle class maps. This repository does not prove the Hodge conjecture, or any case of it.

\begin{itemize}
  \item \textbf{Repository:} \texttt{https://github.com/BenFrohman/uniruled-fourfolds}
  \item \textbf{Zenodo concept DOI (resolves to the latest version):} \texttt{10.5281/zenodo.23652973}
  \item \textbf{Earliest Git Commit SHA:} \texttt{18b62b11b26338ecab62426eaf4844cc57e8353b} (2026-10-01 04:51:29 PM CDT (UTC-5))
  \item \textbf{License \& Copyright:} Copyright \textcopyright\ 2026 Benjamin Stanley Frohman. The text of this paper (LaTeX source and PDF) is licensed under the Creative Commons Attribution 4.0 International License, CC BY 4.0 (see \texttt{LICENSE-paper}). The Lean code is licensed under the Apache License, Version 2.0 (see \texttt{LICENSE} and \texttt{NOTICE}).
\end{itemize}

\section{Setting and named hypotheses}\label{sec:setting}
\begin{definition}[Cycle-class setting, Lean \texttt{CycleClassSetting}]\label{def:setting}
A cycle-class setting $S$ assigns to every fourfold label $X$ types $\mathrm{Hdg}^4(X)$ (rational Hodge classes $H^4(X,\mathbb{Q})\cap H^{2,2}(X)$) and $\mathrm{CH}^2(X)_{\mathbb{Q}}$, and a map $\mathrm{cl}_X:\mathrm{CH}^2(X)_{\mathbb{Q}}\to\mathrm{Hdg}^4(X)$; to every threefold label $Y$ types $\mathrm{Hdg}^2(Y)$ and $\mathrm{CH}^1(Y)_{\mathbb{Q}}$, and a map $\mathrm{cl}_Y:\mathrm{CH}^1(Y)_{\mathbb{Q}}\to\mathrm{Hdg}^2(Y)$; and a predicate ``$X$ is uniruled''. These are uninterpreted data.
\end{definition}

\begin{definition}[Correspondence, Lean \texttt{Correspondence}]\label{def:corr}
A correspondence $\Gamma$ for $X$ consists of a threefold label $\tilde Y$ and maps $\Gamma^*:\mathrm{Hdg}^4(X)\to\mathrm{Hdg}^2(\tilde Y)$, $\Gamma_*:\mathrm{Hdg}^2(\tilde Y)\to\mathrm{Hdg}^4(X)$ and $\Gamma_*:\mathrm{CH}^1(\tilde Y)_{\mathbb{Q}}\to\mathrm{CH}^2(X)_{\mathbb{Q}}$.
\end{definition}

The named hypotheses are:
\begin{itemize}
  \item \textbf{(L11)} \emph{Lefschetz $(1,1)$ on $Y$}, Lean \texttt{Lefschetz11}: for every $\beta\in\mathrm{Hdg}^2(Y)$ there is $z\in\mathrm{CH}^1(Y)_{\mathbb{Q}}$ with $\mathrm{cl}_Y(z)=\beta$.
  \item \textbf{(BS)} \emph{Bloch--Srinivas decomposition on Hodge classes}, Lean \texttt{BSDecomposition}: $\Gamma_*\Gamma^*\gamma=\gamma$ for every $\gamma\in\mathrm{Hdg}^4(X)$.
  \item \textbf{(CL)} \emph{Compatibility of the cycle class map}, Lean \texttt{ClCompatible}: $\mathrm{cl}_X(\Gamma_* z)=\Gamma_*(\mathrm{cl}_{\tilde Y}z)$ for every $z\in\mathrm{CH}^1(\tilde Y)_{\mathbb{Q}}$.
  \item \textbf{(U)} \emph{Bloch--Srinivas correspondences for uniruled fourfolds}, Lean \texttt{BlochSrinivasHypothesis}: every uniruled $X$ has a correspondence $\Gamma$ satisfying (BS) and (CL).
\end{itemize}

\section{Results proved in Lean}\label{sec:results}
Each statement below is checked by the Lean kernel and locked by the Truth Gate (\texttt{TRUTH\_GATE.lock}, \texttt{UniruledFourfolds/Semantics.lean}).

\begin{theorem}[Conditional reduction, Lean \texttt{conditionalCycleClassReduction}]\label{thm:reduction}
Let $S$ be a cycle-class setting, $X$ a fourfold label and $\Gamma$ a correspondence for $X$ with threefold $\tilde Y$. If $\Gamma$ satisfies (BS) and (CL), and (L11) holds on $\tilde Y$, then for every $\gamma\in\mathrm{Hdg}^4(X)$ there is $Z\in\mathrm{CH}^2(X)_{\mathbb{Q}}$ with $\mathrm{cl}_X(Z)=\gamma$.
\end{theorem}
\noindent\emph{Proof.} By (L11) choose $z$ with $\mathrm{cl}_{\tilde Y}(z)=\Gamma^*\gamma$. Then $\mathrm{cl}_X(\Gamma_*z)=\Gamma_*(\mathrm{cl}_{\tilde Y}z)=\Gamma_*\Gamma^*\gamma=\gamma$ by (CL) and (BS). \qed

\begin{theorem}[Conditional statement for uniruled fourfolds, Lean \texttt{uniruled\_cycleClassSurjective}]\label{thm:uniruled}
Let $S$ be a cycle-class setting. If (U) holds and (L11) holds on every threefold label, then for every uniruled fourfold label $X$ and every $\gamma\in\mathrm{Hdg}^4(X)$ there is $Z\in\mathrm{CH}^2(X)_{\mathbb{Q}}$ with $\mathrm{cl}_X(Z)=\gamma$.
\end{theorem}

\begin{proposition}[Consistency, Lean \texttt{hypotheses\_consistent}]\label{prop:consistent}
There are a cycle-class setting $S$, a uniruled fourfold label $X$ and a correspondence $\Gamma$ for $X$ such that (U) holds, (L11) holds on every threefold label, and $\Gamma$ satisfies (BS) and (CL).
\end{proposition}

\begin{proposition}[No hypothesis can be dropped, Lean \texttt{lefschetz\_not\_redundant}, \texttt{decomposition\_not\_redundant}, \texttt{compatibility\_not\_redundant}, \texttt{blochSrinivas\_not\_redundant}]\label{prop:independent}
(a) There are $S$, $X$, $\Gamma$ such that $\Gamma$ satisfies (BS) and (CL), (L11) fails on $\tilde Y$, and some $\gamma\in\mathrm{Hdg}^4(X)$ is not $\mathrm{cl}_X(Z)$ for any $Z$.
(b) There are $S$, $X$, $\Gamma$ such that (L11) holds on $\tilde Y$, $\Gamma$ satisfies (CL), (BS) fails, and some $\gamma\in\mathrm{Hdg}^4(X)$ is not $\mathrm{cl}_X(Z)$ for any $Z$.
(c) There are $S$, $X$, $\Gamma$ such that (L11) holds on $\tilde Y$, $\Gamma$ satisfies (BS), (CL) fails, and some $\gamma\in\mathrm{Hdg}^4(X)$ is not $\mathrm{cl}_X(Z)$ for any $Z$.
(d) There are $S$ and a uniruled $X$ such that (L11) holds on every threefold label, (U) fails, and some $\gamma\in\mathrm{Hdg}^4(X)$ is not $\mathrm{cl}_X(Z)$ for any $Z$.
\end{proposition}

The models in Propositions~\ref{prop:consistent} and~\ref{prop:independent} use the finite types $\mathrm{Unit}$, $\mathrm{Bool}$ and $\mathrm{Empty}$. They show the logical form of the hypotheses: the theorems are not true because their hypotheses are contradictory, and no hypothesis restates the conclusion. They are not geometry.

\section{Open assumptions}\label{sec:open}
The following are assumptions of Theorems~\ref{thm:reduction} and~\ref{thm:uniruled}. None is proved in this repository.
\begin{itemize}
  \item (L11), Lefschetz $(1,1)$: a classical theorem for smooth complex projective varieties (Lefschetz, 1924). Not formalized here.
  \item (BS), the decomposition on Hodge classes, and (CL), compatibility of the cycle class map with the correspondence: modelled on the decomposition of the diagonal of Bloch--Srinivas, Amer.\ J.\ Math.\ \textbf{105} (1983). Not formalized here. This repository does not check that a given fourfold satisfies them in exactly this form.
  \item (U), existence of such correspondences on every uniruled fourfold: the step that Bloch--Srinivas and Conte--Murre take from uniruledness. Not formalized here.
  \item The cohomology groups, Chow groups, cycle class maps and the uniruledness predicate themselves: uninterpreted data (Definition~\ref{def:setting}).
\end{itemize}
In an earlier version, \texttt{Blueprint.lean} postulated these objects with eight \texttt{axiom} declarations and left the reduction as \texttt{sorry}. They are now fields of \texttt{CycleClassSetting} and the named hypotheses above.

\section{Literature, not formalized}
\begin{center}
\small
\begin{tabular}{p{0.5\textwidth}p{0.42\textwidth}}
\toprule
\textbf{Statement} & \textbf{Status} \\
\midrule
Hodge conjecture for a general smooth complex projective variety & Open. \\
Rational Hodge conjecture for fourfolds covered by rational curves & Theorem of Conte--Murre, Math.\ Ann.\ \textbf{238} (1978). Not proved in this repository. \\
Smooth cubic fourfolds & Theorem of Zucker, Compositio Math.\ \textbf{34} (1977). Not proved in this repository. \\
Smooth complex Fano fourfolds & Corollary of Conte--Murre and Koll\'ar--Miyaoka--Mori, J.\ Diff.\ Geom.\ \textbf{36} (1992). Not proved in this repository. \\
Rationality of a general cubic fourfold, Kuznetsov's conjecture & Open. \\
\bottomrule
\end{tabular}
\end{center}
See \texttt{docs/special-cases.md}.

\section{The Mukai lattice file}
\texttt{UniruledFourfolds/MukaiLattice.lean} defines, for integer vectors $v=(v_0,v_2,v_4)$, the form $\langle v,v\rangle=v_2^2-2v_0v_4$ and calls $v$ primitive when $\gcd(|v_0|,|v_2|,|v_4|)=1$, and checks finitely many integer identities about these definitions. This paper states no theorem about it, and it does not bear on cycle-class surjectivity.

\section{Verification}
Continuous integration rebuilds every tracked Lean file (the file list comes from \texttt{git ls-files}; a coverage step names any file that was not built), scans the full build log and every tracked Lean file for \texttt{sorry}, \texttt{admit}, \texttt{native\_decide} and \texttt{axiom} with no exclusions, runs \texttt{\#print axioms} on every theorem, re-checks every module with the external kernel checker, and runs the Truth Gate. The run fails if any theorem depends on an axiom other than \texttt{propext}, \texttt{Classical.choice} and \texttt{Quot.sound}. Rebuild: \texttt{lake build \&\& bash scripts/verify-lock.sh}.

\section{Provenance}
The files \texttt{audit/\allowbreak TITAN\_AUDIT\_CERTIFICATE.json}, \texttt{audit/\allowbreak SEALGUARD\_CLASSIFIED\_MANIFEST.json} and \texttt{PROVENANCE\_\allowbreak MERKLE\_ROOT.json} were committed together at commit \texttt{7f24320} (October 25, 2026). No workflow of this repository produced them, and they are not certificates. As committed, \texttt{TITAN\_AUDIT\_CERTIFICATE.json} recorded \texttt{sorry\_count} $2$ together with the verdict \texttt{PASS\_ZERO\_SORRY\_HONEST\_BOUNDARY}; at that commit \texttt{Blueprint.lean} ended in \texttt{sorry}, and the official GitHub Actions run of that commit reported \texttt{declaration uses `sorry'}, so that verdict was false. The content of all three files is now replaced by pointers to the official \texttt{kernel-certificate} and \texttt{truth-gate-report} artifacts; the originals remain in the git history at \texttt{7f24320}.

The certificates for this paper are the official GitHub Actions artifacts of the \texttt{Kernel certificate} and \texttt{truth-gate} workflows at the commit that contains this paper: the repository's \texttt{lake-manifest.json} and its SHA-256, the toolchain versions, the full \texttt{lake build} log, the \texttt{\#print axioms} output and the external kernel re-check. The file list of any commit is given by \texttt{git ls-files} at that commit.

\vfill
\noindent\rule{\textwidth}{0.4pt}\\
\footnotesize
\textbf{Attribution \& Licensing Notice:}\\
Copyright \textcopyright\ 2026 Benjamin Stanley Frohman (ORCID: \texttt{0009-0006-7068-3718}).\\
Paper text: CC BY 4.0 (\texttt{LICENSE-paper}). Lean code: Apache License, Version 2.0 (\texttt{LICENSE}). Reuse must credit Benjamin Stanley Frohman and keep the \texttt{NOTICE}.
\end{document}
